% ===== ABSTRACT (200-350 words, ONE flowing text, no bullet points) =====
\addcontentsline{toc}{chapter}{ABSTRACT}
\begin{center}\textbf{\large ABSTRACT}\end{center}

% GUIDE: write 4 parts as continuous prose:
%  (i)   the problem: small/medium distributors still run purchasing, warehousing,
%        sales and delivery on paper/Excel -> stockouts, expired goods, no traceability.
%  (ii)  chosen approach: a single web-based supply-chain management system with a
%        layered Spring Boot + React architecture and role-based workflows.
%  (iii) solution overview: integrated procure-to-stock, order-to-cash and delivery
%        modules with FEFO lot tracking and inventory reservation.
%  (iv)  main contribution & results: FEFO allocation algorithm, reservation logic,
%        failed-delivery recovery; summarise concrete results (modules, screens, tests).

\noindent Small and medium-sized distributors often manage purchasing,
warehousing, sales, invoicing, and delivery through disconnected spreadsheets,
paper documents, and manual communication. This practice leads to inaccurate
stock information, overselling, slow approval, weak traceability, and a high
risk of expired goods remaining in the warehouse. This thesis designs and
implements a web-based distribution-management system that unifies these
activities behind role-based workflows and lot-level inventory control. The
system follows a layered architecture with a Spring Boot REST backend, a React
single-page frontend, and a PostgreSQL database, integrating procure-to-stock,
order-to-cash, inventory, invoicing, accounting, delivery, reporting, and
administration modules in one application. It maintains inventory at both the
aggregate level and the lot level, where each lot carries a batch number and an
expiry date: goods receipt creates inventory lots, sales-order approval reserves
available stock to prevent overselling, and goods issue deducts stock following
the First-Expired-First-Out (FEFO) principle. A recovery mechanism for failed
delivery restores reserved inventory, adjusts the delivered quantity, updates
the order status, and corrects the related invoice. Role-Based Access Control is
enforced for nine business roles, separating document creation, approval,
warehouse confirmation, delivery assignment, shipper execution, and
administration. The implemented prototype comprises a PostgreSQL schema of 32
tables, a Spring Boot backend of 21 REST controllers exposing 182 endpoints
across nine packages (176 Java classes, 11,847 lines of code), and a React
frontend of 39 screens (10,923 lines), all measured with cloc excluding blank
and comment lines. In a representative performance evaluation, removing an N+1
query pattern in the available-delivery listing reduced the endpoint response
time from 2,380 ms to 47 ms. The main contribution of the thesis is the
integration of FEFO lot allocation, inventory reservation, failed-delivery
recovery, and role-based workflow enforcement into a single, focused
distribution-management system.

\vspace{0.8cm}
\noindent Student

\hfill \textit{Duong Phuong Thao}
